%! Polar Function Graphs — Unit 7, Lesson 7.7 — Guided Notes
\documentclass[10pt]{article}
\usepackage{precalculus-article}
\usepackage{precalculus-boxes}
\newcommand{\TallMath}[1]{$\displaystyle #1\rule[-1.4em]{0pt}{3.2em}$}

\begin{document}

\pageheader{Unit 7, Lesson 7.7}{Guided Notes}
\namedateperiod

\vspace{0.08in}

%----------------------------------------------------------------------
% Objectives
%----------------------------------------------------------------------
\begin{objectivebox}
\textbf{I will be able to\ldots}
\begin{enumerate}[leftmargin=1.5em, itemsep=4pt]
  \item \textbf{LO 3.14.A} --- Read a polar function $r = f(\theta)$ from a table of
    input-output pairs and describe how the radius changes as the angle increases.
  \item \textbf{LO 3.14.A} --- Connect a rectangular graph of $r = f(\theta)$ to the
    corresponding polar curve by identifying zeros, maxima, and minima.
  \item \textbf{LO 3.14.A} --- Describe which portion of a polar curve is traced when
    the domain of $\theta$ is restricted to a sub-interval.
\end{enumerate}
\end{objectivebox}

\vspace{0.08in}

%----------------------------------------------------------------------
% Vocabulary
%----------------------------------------------------------------------
\begin{vocabbox}
\begin{description}[leftmargin=0pt, itemsep=6pt]
  \item[\termblanklong{polar function}]
    A function $r = f(\theta)$ where $\theta$ (the \underline{\hspace{2cm}})
    is the input and $r$ (the \underline{\hspace{2cm}}) is the output.

  \item[\termblanklong{pole}]
    The \underline{\hspace{2.5cm}} in polar coordinates. The curve passes through
    the pole whenever $r = $ \underline{\hspace{1cm}}.

  \item[\termblanklong{limaçon}]
    A polar curve of the form $r = a + b\sin\theta$ or $r = a + b\cos\theta$.
    When $|a| = |b|$ it is called a \underline{\hspace{2cm}}.

  \item[\termblanklong{rose curve}]
    A polar curve of the form $r = a\sin(n\theta)$ or $r = a\cos(n\theta)$.
    It has $n$ petals when $n$ is odd and \underline{\hspace{1cm}} petals when
    $n$ is even.

  \item[\termblanklong{petal}]
    One lobe of a rose curve; traced as $r$ goes from \underline{\hspace{0.8cm}}
    to a maximum and back to \underline{\hspace{0.8cm}}.
\end{description}
\end{vocabbox}

\vspace{0.08in}

%----------------------------------------------------------------------
% Hook
%----------------------------------------------------------------------
\begin{hookbox}
\textbf{Entry Question:} The graph below shows the polar curve $r = 2\sin\theta$.
It turns out to be a \emph{circle}.

\vspace{4pt}
\begin{center}
\begin{tikzpicture}[x=1.4cm, y=1.4cm]
  % Polar grid (light circles and rays)
  \foreach \r in {0.5, 1, 1.5, 2} {
    \draw[linegray, very thin] (0,0) circle (\r);
  }
  \foreach \ang in {0, 30, 60, 90, 120, 150, 180, 210, 240, 270, 300, 330} {
    \draw[linegray, very thin] (0,0) -- ({2.1*cos(\ang)}, {2.1*sin(\ang)});
  }
  % Axes
  \draw[->] (-2.2,0) -- (2.3,0) node[right]{\small $x$};
  \draw[->] (0,-2.2) -- (0,2.3) node[above]{\small $y$};
  % r = 2sin(theta) in parametric form: x = 2sin(t)cos(t), y = 2sin^2(t)
  \draw[navy, very thick, domain=0:3.1416, samples=200, smooth]
    plot ({2*sin(deg(\x))*cos(\x r)}, {2*sin(deg(\x))*sin(\x r)});
  % Labels
  \node[right, font=\small, navy] at (0.1, 1.1) {$r = 2\sin\theta$};
  \node[below right, font=\scriptsize] at (1,0) {$1$};
  \node[below right, font=\scriptsize] at (2,0) {$2$};
  \node[above right, font=\scriptsize] at (0,1) {$1$};
  \node[above right, font=\scriptsize] at (0,2) {$2$};
\end{tikzpicture}
\end{center}

\vspace{4pt}
\textbf{Q1.} The maximum value of $r$ on this curve is \underline{\hspace{1.5cm}},
occurring at $\theta = $ \underline{\hspace{1.5cm}}.

\vspace{4pt}
\textbf{Q2.} The value of $r$ when $\theta = 0$ is \underline{\hspace{1.5cm}}.
This means the curve \underline{\hspace{3cm}} (passes through / does not pass through)
the pole at $\theta = 0$.

\vspace{4pt}
\textbf{Q3.} How could a sine function produce a circle? \writeline

\end{hookbox}

\vspace{0.08in}

%----------------------------------------------------------------------
% LO 3.14.A — Part 1: Reading from a Table
%----------------------------------------------------------------------
\begin{notesbox}{LO 3.14.A --- Reading Polar Functions from Tables}

The table below gives input-output pairs for $r = 2\sin\theta$.

\vspace{4pt}
\renewcommand{\arraystretch}{1.5}
\begin{tabular}{|c|c|c|c|c|c|c|c|}
\hline
$\theta$ & $0$ & $\dfrac{\pi}{6}$ & $\dfrac{\pi}{3}$ & $\dfrac{\pi}{2}$ &
           $\dfrac{2\pi}{3}$ & $\dfrac{5\pi}{6}$ & $\pi$ \\
\hline
$r = 2\sin\theta$ & $0$ & $1$ & $\sqrt{3} \approx 1.73$ & $2$ &
                   $\sqrt{3} \approx 1.73$ & $1$ & $0$ \\
\hline
\end{tabular}

\vspace{8pt}
Use the table to complete the sentences:

\vspace{4pt}
As $\theta$ increases from $0$ to $\dfrac{\pi}{2}$, $r$
\underline{\hspace{2.5cm}} (increases / decreases) from
\underline{\hspace{0.8cm}} to \underline{\hspace{0.8cm}}.

\vspace{4pt}
As $\theta$ increases from $\dfrac{\pi}{2}$ to $\pi$, $r$
\underline{\hspace{2.5cm}} (increases / decreases) from
\underline{\hspace{0.8cm}} to \underline{\hspace{0.8cm}}.

\vspace{4pt}
The curve passes through the pole (where $r = 0$) at $\theta = $
\underline{\hspace{1cm}} and $\theta = $ \underline{\hspace{1cm}}.

\vspace{4pt}
The maximum distance from the origin is $r = $ \underline{\hspace{1cm}},
occurring at $\theta = $ \underline{\hspace{1.5cm}}.

\end{notesbox}

\vspace{0.08in}

%----------------------------------------------------------------------
% LO 3.14.A — Part 2: Connecting Rectangular and Polar Graphs
%----------------------------------------------------------------------
\begin{notesbox}{LO 3.14.A --- Matching Rectangular and Polar Graphs}

The rectangular graph of $r = \sin(2\theta)$ is shown below for $\theta \in [0, 2\pi]$.

\vspace{4pt}
\begin{center}
\begin{tikzpicture}[x=1.1cm, y=1.2cm]
  \draw[linegray, very thin] (-0.2,-1.3) grid[xstep=0.5, ystep=0.5] (6.7,1.3);
  \draw[->] (-0.2,0) -- (6.9,0) node[right]{\small $\theta$};
  \draw[->] (0,-1.35) -- (0,1.4) node[above]{\small $r$};
  \foreach \x/\lbl in {0.785/$\frac{\pi}{4}$, 1.571/$\frac{\pi}{2}$, 2.356/$\frac{3\pi}{4}$,
                        3.142/$\pi$, 3.927/$\frac{5\pi}{4}$, 4.712/$\frac{3\pi}{2}$,
                        5.498/$\frac{7\pi}{4}$, 6.283/$2\pi$} {
    \draw (\x,0.06) -- (\x,-0.06);
    \node[below, font=\scriptsize] at (\x,-0.06) {\lbl};
  }
  \foreach \y/\lbl in {-1/{$-1$}, 1/{$1$}} {
    \draw (0.07,\y) -- (-0.07,\y);
    \node[left, font=\scriptsize] at (-0.07,\y) {\lbl};
  }
  \draw[navy, thick, domain=0:6.3, samples=200]
    plot (\x, {sin(2*deg(\x))});
  \node[right, font=\small, navy] at (6.4, 0.3) {$r = \sin(2\theta)$};
\end{tikzpicture}
\end{center}

\vspace{4pt}
\textbf{Reading the rectangular graph:}

The value of $r = 0$ at $\theta = $ \underline{\hspace{1cm}},
\underline{\hspace{1cm}}, \underline{\hspace{1cm}}, \underline{\hspace{1cm}},
\underline{\hspace{1cm}}.
Each zero corresponds to the polar curve \underline{\hspace{3.5cm}} (passing through / leaving)
the pole.

\vspace{4pt}
The maximum value of $r$ is \underline{\hspace{0.8cm}}, occurring at
$\theta = $ \underline{\hspace{1.5cm}} and $\theta = $ \underline{\hspace{1.5cm}}.
Each maximum corresponds to the \underline{\hspace{2cm}} (tip / base) of a petal.

\vspace{4pt}
The number of zeros between $\theta = 0$ and $\theta = 2\pi$ (not counting $\theta = 0$)
is \underline{\hspace{0.8cm}}. The polar curve has \underline{\hspace{0.8cm}} petals.

\end{notesbox}

\vspace{0.08in}

%----------------------------------------------------------------------
% LO 3.14.A — Part 3: Domain Restriction
%----------------------------------------------------------------------
\begin{notesbox}{LO 3.14.A --- Restricting the Domain}

\textbf{Key idea (EK 3.14.A.2):} If we restrict $\theta$ to a sub-interval, we trace
only \emph{part} of the polar curve.

\vspace{6pt}
Using the rectangular graph of $r = \sin(2\theta)$ above, answer the following:

\vspace{4pt}
\textbf{1.}\quad When $\theta \in \left[0, \dfrac{\pi}{2}\right]$:
$r$ goes from \underline{\hspace{0.8cm}} up to \underline{\hspace{0.8cm}} and back
to \underline{\hspace{0.8cm}}.
This traces \underline{\hspace{0.8cm}} petal(s) of the polar curve.

\vspace{4pt}
\textbf{2.}\quad When $\theta \in [0, \pi]$:
This restriction traces \underline{\hspace{0.8cm}} petal(s) of the polar curve.
Explain: \writeline

\vspace{4pt}
\textbf{3.}\quad When $\theta \in [0, 2\pi]$:
This traces \underline{\hspace{0.8cm}} petal(s) total.

\end{notesbox}

\vspace{0.08in}

%----------------------------------------------------------------------
% Practice
%----------------------------------------------------------------------
\begin{practicebox}

\textbf{Guided Practice 1.}\quad The table below gives values for $r = 1 + \cos\theta$.
Use the table to answer the questions.

\vspace{4pt}
\renewcommand{\arraystretch}{1.5}
\begin{tabular}{|c|c|c|c|c|c|c|}
\hline
$\theta$ & $0$ & $\dfrac{\pi}{3}$ & $\dfrac{\pi}{2}$ & $\pi$ & $\dfrac{3\pi}{2}$ & $2\pi$ \\
\hline
$r = 1 + \cos\theta$ & $2$ & $1.5$ & $1$ & $0$ & $1$ & $2$ \\
\hline
\end{tabular}

\vspace{6pt}
(a)\quad The maximum value of $r$ is \underline{\hspace{1cm}}, at $\theta = $ \underline{\hspace{1.5cm}}.

\vspace{4pt}
(b)\quad The curve passes through the pole at $\theta = $ \underline{\hspace{1.5cm}}.

\vspace{0.10in}
\textbf{Guided Practice 2.}\quad For $r = \sin(2\theta)$, the maximum value of $r$ is
\underline{\hspace{0.8cm}}. It occurs at $\theta = \dfrac{\pi}{4}$.
What does this tell you about the tip of the first petal? \writeline

\end{practicebox}

\end{document}
